\section{An exact negative certificate for $\mathbf L_\star$}
\label{app:negative-witness}

This appendix certifies the only finite computation used in
\Cref{lem:negative-limit}: the fixed matrix $\mathbf L_\star$ has a
strictly negative eigenvalue.  The construction uses exactly the two
vertical families of test speeds listed in the main text.

For the certificate only, it is convenient to parameterize the same
finite set by
\[
 I=\{(r,j,\sigma):r\in\{1,2\},\ 1\le j\le256,\
 \sigma\in\{-1,1\}\},
\]
with
\[
 a_{r,j,\sigma}=r+i\sigma j.
\]
For $\ell=(r,j,\sigma)$ and $h=(r',j',\sigma')$, the entries of
$\mathbf L_\star$ are
\begin{equation}\label{C:eq:entry}
 (\mathbf L_\star)_{\ell h}
 =
 2\left(
 \frac{2}{A_{\ell h}^2}
 +\frac{1}{C_{\ell h}^2}
 -\frac{1}{A_{\ell h}C_{\ell h}}
 \right),
\end{equation}
where
\[
 A_{\ell h}
 =(r+r')^2+(\sigma j-\sigma'j')^2,\qquad
 C_{\ell h}=4rr'.
\]
Thus every entry is rational.

Define the integer vector $d\in\mathbb Z^{1024}$ by
\begin{equation}\label{C:eq:witness}
 d_{r,j,\sigma}
 :=
 \sigma\widehat d_j,\qquad
 \widehat d_j
 :=
 \left\lfloor
 \frac{2^{43}j}{(128^2+j^2)^2}
 \right\rfloor.
\end{equation}

\begin{lemma}[Exact negative direction]\label{C:lem:witness}
The vector in \eqref{C:eq:witness} satisfies
\begin{equation}\label{C:eq:witness-bounds}
 \|d\|_1=859652168<2^{30},
 \qquad
 d^\top\mathbf L_\star d<-2^{42}.
\end{equation}
Consequently,
\begin{equation}\label{C:eq:eig-margin}
 \lambda_{\min}(\mathbf L_\star)<-2^{-18}.
\end{equation}
In particular, the constant $\gamma$ in
\Cref{lem:negative-limit} may be chosen as $\gamma=2^{-18}$.
\end{lemma}

\begin{proof}
Since the entries in \eqref{C:eq:entry} are rational, the quadratic
form can be checked using integer arithmetic only.  Define
\[
 \mathcal J
 :=
 \sum_{\ell,h\in I}
 \left\lfloor
 \frac{2^{64}d_\ell d_h
 \bigl(4C_{\ell h}^2+2A_{\ell h}^2
       -2A_{\ell h}C_{\ell h}\bigr)}
 {A_{\ell h}^2C_{\ell h}^2}
 \right\rfloor.
\]
Each summand is rounded toward $-\infty$.  Since $|I|^2=2^{20}$,
\[
 \frac{\mathcal J}{2^{64}}
 \le d^\top\mathbf L_\star d
 \le \frac{\mathcal J+2^{20}}{2^{64}}.
\]
The exact integer evaluation is
\[
\begin{aligned}
 \mathcal J
 &=-111807741093765987869670624504700,\\
 \mathcal J+2^{20}
 &=-111807741093765987869670623456124,
\end{aligned}
\]
whose upper numerator is smaller than $-2^{106}$.  Hence
$d^\top\mathbf L_\star d<-2^{42}$.  The stated $\ell_1$ norm follows
directly from \eqref{C:eq:witness}.

Now
\[
 \|d\|_2^2\le\|d\|_1^2<2^{60}.
\]
Therefore the Rayleigh quotient satisfies
\[
 \frac{d^\top\mathbf L_\star d}{\|d\|_2^2}
 <-\frac{2^{42}}{2^{60}}
 =-2^{-18}.
\]
By the variational characterization of the smallest eigenvalue,
\[
 \lambda_{\min}(\mathbf L_\star)
 \le
 \frac{d^\top\mathbf L_\star d}{\|d\|_2^2}
 <-2^{-18},
\]
which proves \eqref{C:eq:eig-margin}.
\end{proof}

\begin{remark}
A direct floating-point diagonalization gives
$
\lambda_{\min}(\mathbf L_\star)\approx-1.18428\times10^{-2},
$
which is substantially more negative than the conservative exact
margin in \eqref{C:eq:eig-margin}.  The numerical value is not used in
the proof.
\end{remark}

\begin{remark}[Alternative families of test points]
\label{C:rem:alternative-test-families}
The two-family choice $\Re a=1,2$ is not unique.  Preliminary numerical
experiments with three vertical families
\[
 \Re a\in\{1,2,3\},\qquad
 \Im a=\pm1,\ldots,\pm J,
\]
show a negative eigenvalue already near $J=44$, corresponding to
$s=6J=264$ test points.  We retain the two-family construction because
it is the family for which the exact certificate above is supplied.
\end{remark}
